\chapter{MACSLAM: Unified Uncertainty-Aware Multi-Frame Visual-Inertial SLAM} \label{chapter:macslam}
\fancyhead[LO]{\makebox[\textwidth][r]{Chapter 9. MACSLAM}}

\section{Introduction}

This chapter presents MACSLAM, a unified SLAM framework that combines multi-frame optimization, visual-inertial fusion, and loop closure under a shared uncertainty-aware formulation.
It extends modules from earlier chapters: metrics-aware visual covariance from MAC-VO, inertial uncertainty modeling from AirIMU and AirIO, and adaptive visual-inertial fusion from MACVIO (Chapter~\ref{chapter:macvio}).

The key design choice is a unified visual foundation front-end based on DINO features.
The same visual representation is used for local geometric matching, global visual place recognition, and semantic-language alignment through lightweight adapters.
This design aims to reduce front-end fragmentation and make local odometry, loop closure, and semantic mapping mutually consistent.

The motivation for MACSLAM is that robust autonomy requires more than local odometry.
A deployable system must maintain global consistency over long trajectories, recover from drift through loop closure, and support higher-level map understanding for planning and interaction.
Existing pipelines often treat these capabilities as loosely connected modules with separate feature spaces, which introduces inconsistency and brittle behavior under domain shift.

MACSLAM addresses this fragmentation through two unification principles.
First, visual and inertial information are fused in one uncertainty-aware multi-frame optimizer rather than isolated local modules.
Second, one foundation visual representation supports local matching, global place recognition, and semantic-language mapping.
This shared representation is intended to improve both robustness and scalability while opening a path toward language-aware semantic SLAM in real deployments.

\section{Method}
\label{sec:macslam:method}

\subsection{Unified Optimization Objective}

Let $\mathcal{X}_{1:T}$ be the trajectory states, $\mathcal{M}$ the map variables, and $\mathcal{L}$ loop-closure constraints.
MACSLAM solves a single factor-graph objective:
\begin{equation}
\label{eq:macslam:objective}
(\mathcal{X}_{1:T}^{*}, \mathcal{M}^{*})
=
\argmin_{\mathcal{X}_{1:T},\mathcal{M}}
\sum_{k\in\mathcal{V}} \left\| r_k^{v} \right\|_{{\Sigma}_k^{v}}^{2}
+
\sum_{k\in\mathcal{I}} \left\| r_k^{i} \right\|_{{\Sigma}_k^{i}}^{2}
+
\sum_{k\in\mathcal{L}} \left\| r_k^{lc} \right\|_{{\Sigma}_k^{lc}}^{2},
\end{equation}
where visual, inertial, and loop-closure residuals are jointly optimized with learned, metrics-aware covariance.
Unlike conventional pipelines with fixed noise assumptions, MACSLAM propagates data-conditioned uncertainty through all factor types.

\subsection{DINO-Based Unified Visual Front-End}

MACSLAM uses DINO features as a shared visual backbone with three heads:
\begin{itemize}
    \item \textbf{Geometric head:} dense feature matching and correspondence confidence for local visual odometry factors.
    \item \textbf{Global retrieval head:} place descriptors for loop-closure candidate retrieval and verification.
    \item \textbf{Semantic-language head:} adapter from DINO tokens to semantic and language-aligned embeddings for semantic mapping.
\end{itemize}

The core hypothesis is that one foundation representation can support both local geometry and global semantics more consistently than separate task-specific front-ends.
This allows loop closure and semantic mapping to directly reuse features already optimized for odometry.

\subsection{Visual-Inertial Fusion in Multi-Frame SLAM}

Visual and inertial factors are fused under uncertainty-aware weighting.
Visual covariance priors come from the metrics-aware modeling of MAC-VO, while inertial covariance priors come from AirIMU/AirIO and are adapted online via MACVIO-style reliability control.
Fusion is not implemented as static modality weighting; instead, confidence is adjusted online according to degradation context (for example, motion blur, low texture, aggressive maneuvers, and IMU disturbance).

\subsection{Loop Closure and Semantic Mapping}

Loop closure is modeled as a two-stage process: DINO-based global retrieval proposes candidates, and geometric verification adds loop factors with uncertainty-aware covariance.
Candidates with high uncertainty are down-weighted to reduce false closures.
In parallel, semantic-language adapters project visual tokens into language-aware map attributes, enabling semantic mapping that can support text-conditioned retrieval and higher-level scene reasoning.

\subsection{Summary}

MACSLAM provides:
\begin{itemize}
    \item improved trajectory consistency through unified multi-frame visual-inertial optimization,
    \item more reliable loop closure under domain shift via uncertainty-aware global matching,
    \item reduced system complexity by using one DINO-based front-end for geometry, retrieval, and semantics,
    \item semantic mapping with language-aligned features on top of the same SLAM backbone.
\end{itemize}
These capabilities support the thesis goal of a full spatial perception system that is accurate, robust, and deployable with minimal additional human engineering.